\section{引言}

量子色动力学（QCD）最重要的目标之一，是从其基本自由度——夸克和胶子——出发理解强子结构。格点 QCD 是获得这类结果的一个有力工具，它已被用于较精确地研究强子的静态性质，例如质量、电荷半径等（参见文献~\cite{Edwards:2012fx,Aoki:2016frl}）。然而一般而言，它无法直接处理具有实时依赖性的物理量。一个例子是部分子分布函数（PDF）：其定义为光锥关联的夸克与胶子矩阵元，描述强子内部夸克和胶子的动量分布。它们对于理解诸如 LHC 这类高能强子对撞机的实验数据起着关键作用。格点 QCD 只能通过计算矩来间接提取 PDF 的信息~\cite{Detmold:2001jb,Dolgov:2002zm}，而这受到技术困难的限制。

过去几年中，人们提出了一种从格点 QCD 出发研究部分子物理的新方法，即如今所称的大动量有效理论（LaMET）~\cite{Ji:2013dva,Ji:2014gla}。在这一方法中，人们从一个与时间无关的准分布（quasi-PDF）出发——它是可以直接在格点上计算的非局域夸克与胶子算符的矩阵元——然后借助微扰因子化用它来抽取光锥 PDF，误差为被核子动量压低的幂次修正。文献~\cite{Ji:2013dva,Xiong:2013bka} 中给出的因子化是针对裸夸克准分布的，即进入准分布的所有场和耦合常数都是裸量。
准 PDF 的紫外（UV）物理全部被微扰地因子化到匹配系数之中。然而，众所周知格点微扰论收敛缓慢，通常需要非微扰重正化来定义连续极限，以便更好地控制与部分子物理的匹配。因此，要让 LaMET 更加实用，就必须完整理解相关非局域算符的 UV 发散。

QCD 中非局域算符的重正化在文献中尚未得到充分研究。LaMET 非局域算符的一个显著特征，是出现了连接不同时空点处部分子场的威尔逊线以保证规范不变性。
关于开放威尔逊线以及闭合威尔逊圈的重正化，早在几十年前就已被研究~\cite{Dotsenko:1979wb,Craigie:1980qs,Dorn:1986dt}，结果表明具有光滑简单回路的威尔逊线可以被重正化，特别是与威尔逊线相联系的幂次发散可以被吸收进指数形式的质量重正化之中。在辅助场表述~\cite{Dorn:1986dt} 中，还证明了威尔逊线可以替换为辅助场的格林函数并据此加以研究。另一方面，诸如准 PDF 这类非局域算符的重正化，由于端点处连有部分子场而具有额外的复杂性。
我们此前已证明，准 PDF 在维数正规化下直到两圈阶都是可乘法重正化的~\cite{Ji:2015jwa}。
准 PDF 完整的重正化性质在近期的若干研究中已被猜想~\cite{Ishikawa:2016znu,Chen:2016fxx,Constantinou:2017sej,Alexandrou:2017huk,Chen:2017mzz,Stewart:2017ss}（关于对横动量依赖分布重正化的类似猜想的更早工作，见~\cite{Musch:2010ka}），
但尚未被证明。

在本文中，我们对准 PDF 的规范不变非局域算符的重正化进行系统研究。我们首先给出一个一般性框架：在 QCD 拉氏量中引入辅助的“重夸克”场，于是威尔逊线可以替换为辅助场的乘积。通过与重夸克有效理论（HQET）类比，我们论证该理论在微扰论所有阶都可重正化。随后以无极化夸克准 PDF 为例，我们证明其重正化归结为两个轻重夸克流的重正化，后者同样可以在微扰论所有阶完成。这消除了通过格点模拟定义连续准 PDF 的障碍。同一结论并不局限于夸克准 PDF，还可以推广到胶子准 PDF，以及其他非局域关联函数，例如 LaMET 框架中定义的、用以抽取广义部分子分布、横动量依赖分布、强子分布振幅等的欧几里得观测量。
